\documentclass[a4paper]{article}

\usepackage[spanish]{babel}
\usepackage{graphicx}
\usepackage{amsmath, amssymb, mathrsfs}
\usepackage[margin=2cm]{geometry}
\usepackage{fancyhdr}
\usepackage{enumerate}
\usepackage[shortlabels]{enumitem}
\usepackage{parskip}
\usepackage[most]{tcolorbox}
\usepackage[hidelinks]{hyperref}
\usepackage{float}
\usepackage{booktabs}
\usepackage{mathrsfs}
\usepackage{tikz}
\usetikzlibrary{arrows.meta, calc, decorations.pathreplacing}
\usepackage{pgfplots}
\pgfplotsset{compat=1.18}


% cabecera
\pagestyle{fancy}
\fancyhead[l]{Física Matemática}
\fancyhead[c]{Parcial 3}
\fancyhead[r]{\today}
\fancyfoot[c]{\thepage}
\renewcommand{\headrulewidth}{0.2pt}


\begin{document}


\section{Problema 3}
Considere un cascarón esférico de radio $R$ cuyo potencial electrostático sobre la superficie es
$$\phi(R,\theta) = V_0 \sin^2\!\left(\frac{\theta}{2}\right)$$
\begin{enumerate}[a)]
\item Determinar el potencial eléctrico $\phi(r,\theta)$ en el interior de la esfera.
\item Hallar las componentes radial $E_r$ y transversal $E_\theta$ del campo eléctrico en el interior de la esfera.
\end{enumerate}

\subsection{Solución Problema 3 — Método de Ortogonalidad}

\textbf{item a)}

Para el interior de la esfera ($r < R$), debemos imponer $B_l = 0$ para garantizar la regularidad
en $r = 0$. La solución general queda:

$$\phi(r,\theta) = \sum_{l=0}^{\infty} A_l\, r^l\, P_l(\cos\theta)$$

Evaluando en $r = R$ e igualando a la condición de frontera:
$$\sum_{l=0}^{\infty} A_l R^l\, P_l(\cos\theta) = V_0 \sin^2\!\left(\frac{\theta}{2}\right)$$

Para hallar los coeficientes $A_l$ usamos \textbf{explícitamente la condición de ortogonalidad} de
los polinomios de Legendre. Multiplicando ambos lados por $P_{l'}(\cos\theta)\sin\theta$ e
integrando en $[0,\pi]$, y usando que
$$\int_{-1}^{1} P_l(u)\,P_{l'}(u)\,du = \frac{2}{2l+1}\,\delta_{l\,l'},$$
obtenemos directamente:
\begin{equation}
    A_l R^l = \frac{2l+1}{2}\int_0^{\pi} V_0\sin^2\!\left(\frac{\theta}{2}\right)
    P_l(\cos\theta)\,\sin\theta\, d\theta
    \label{eq:coef_Al}
\end{equation}

\textbf{Evaluación de la integral sin convertir $\sin$ a $\cos$ a priori.}
Realizamos la sustitución estándar para integrales de Legendre: $u = \cos\theta$,
$du = -\sin\theta\,d\theta$. Al hacer este cambio de variable, la identidad trigonométrica
$\sin^2(\theta/2) = (1-\cos\theta)/2$ surge de forma \emph{natural} como paso algebraico
ordinario —no como observación estratégica previa— dando:
$$\sin^2\!\left(\frac{\theta}{2}\right) = \frac{1-u}{2}$$

El límite de integración cambia de $[0,\pi]$ a $[1,-1]$, absorbiendo el signo del jacobiano:
$$A_l R^l = \frac{(2l+1)V_0}{2}\int_{-1}^{1} \frac{1-u}{2}\,P_l(u)\,du
           = \frac{(2l+1)V_0}{4}
             \left[\int_{-1}^{1} P_l(u)\,du - \int_{-1}^{1} u\,P_l(u)\,du\right]$$

Evaluamos cada integral por ortogonalidad, reconociendo que $1 = P_0(u)$ y $u = P_1(u)$:

\begin{align}
    \int_{-1}^{1} P_l(u)\,du &= \int_{-1}^{1} P_0(u)\,P_l(u)\,du
        = \frac{2}{2\cdot 0+1}\,\delta_{l,0} = 2\,\delta_{l,0}
    \label{eq:int1}\\[4pt]
    \int_{-1}^{1} u\,P_l(u)\,du &= \int_{-1}^{1} P_1(u)\,P_l(u)\,du
        = \frac{2}{2\cdot 1+1}\,\delta_{l,1} = \frac{2}{3}\,\delta_{l,1}
    \label{eq:int2}
\end{align}

Sustituyendo \eqref{eq:int1} y \eqref{eq:int2} en \eqref{eq:coef_Al}:
$$A_l R^l = \frac{(2l+1)V_0}{4}\left[2\,\delta_{l,0} - \frac{2}{3}\,\delta_{l,1}\right]$$

Calculamos cada coeficiente:
\begin{itemize}
    \item $l = 0$: $\displaystyle A_0 = \frac{1\cdot V_0}{4}\cdot 2 = \frac{V_0}{2}$
    \item $l = 1$: $\displaystyle A_1 R = \frac{3V_0}{4}\cdot\left(-\frac{2}{3}\right) = -\frac{V_0}{2}$,
          luego $A_1 = -\dfrac{V_0}{2R}$
    \item $l \ge 2$: $A_l = 0$
\end{itemize}

Aplicando los tres puntos anteriores en la expansión general:
$$\phi(r,\theta) = \sum_{l=0}^{\infty} A_l r^l P_l(\cos\theta)
= \frac{V_0}{2}P_0(\cos\theta) - \frac{V_0}{2R}\,r\,P_1(\cos\theta)
+ \underbrace{\sum_{l=2}^{\infty} 0}_{=\,0}$$

\begin{tcolorbox}
$$\phi(r,\theta) = \frac{V_0}{2} - \frac{V_0}{2R}\,r\cos\theta$$
\end{tcolorbox}

\textbf{item b)}

El campo eléctrico se obtiene de $\vec{E} = -\vec{\nabla}\phi$.
Con simetría azimutal el gradiente en esféricas es:
$$\vec{\nabla}\phi = \frac{\partial\phi}{\partial r}\,\hat{e}_r
                  + \frac{1}{r}\frac{\partial\phi}{\partial\theta}\,\hat{e}_\theta$$

Componente radial:
$$E_r = -\frac{\partial\phi}{\partial r} = -\!\left(-\frac{V_0}{2R}\cos\theta\right)
      = \frac{V_0}{2R}\cos\theta$$

Componente transversal:
$$E_\theta = -\frac{1}{r}\frac{\partial\phi}{\partial\theta}
           = -\frac{1}{r}\left(\frac{V_0}{2R}\,r\sin\theta\right)
           = -\frac{V_0}{2R}\sin\theta$$

El campo resultante es uniforme (en la dirección $\hat{e}_k \equiv \hat{z}$):
$$\vec{E} = \frac{V_0}{2R}\!\left(\cos\theta\,\hat{e}_r - \sin\theta\,\hat{e}_\theta\right)
= \frac{V_0}{2R}\,\hat{e}_k$$

\subsubsection*{item c)}

Para la región exterior ($r > R$), imponemos $A_l = 0$ para que $\phi$ sea finito cuando
$r\to\infty$:
\begin{equation}
    \phi(r,\theta) = \sum_{l=0}^{\infty} \frac{B_l}{r^{l+1}}\,P_l(\cos\theta)
\end{equation}

Evaluando en $r = R$ e igualando a la condición de frontera, la condición de ortogonalidad entrega:
\begin{equation}
    \frac{B_l}{R^{l+1}} = \frac{2l+1}{2}\int_0^{\pi} V_0\sin^2\!\left(\frac{\theta}{2}\right)
    P_l(\cos\theta)\,\sin\theta\,d\theta
    \label{eq:coef_Bl}
\end{equation}

El lado derecho de \eqref{eq:coef_Bl} es \emph{exactamente el mismo integral} que en
\eqref{eq:coef_Al}, por lo que no hace falta repetir el cálculo:
$$\frac{B_l}{R^{l+1}} = \frac{(2l+1)V_0}{4}\left[2\,\delta_{l,0} - \frac{2}{3}\,\delta_{l,1}\right]$$

Calculamos cada coeficiente:
\begin{itemize}
    \item $l = 0$: $\dfrac{B_0}{R} = \dfrac{V_0}{2}$, luego $B_0 = \dfrac{V_0 R}{2}$
    \item $l = 1$: $\dfrac{B_1}{R^2} = -\dfrac{V_0}{2}$, luego $B_1 = -\dfrac{V_0 R^2}{2}$
    \item $l \ge 2$: $B_l = 0$
\end{itemize}

\begin{tcolorbox}
$$\phi(r,\theta) = \frac{V_0 R}{2r} - \frac{V_0 R^2}{2r^2}\cos\theta$$
\end{tcolorbox}

\subsubsection*{item d)}

Aplicando $\vec{E} = -\vec{\nabla}\phi$ al potencial exterior:

Componente radial:
\begin{equation}
    E_r = -\frac{\partial\phi}{\partial r}
        = -\!\left(-\frac{V_0 R}{2r^2} + \frac{V_0 R^2}{r^3}\cos\theta\right)
        = \frac{V_0 R}{2r^2} - \frac{V_0 R^2}{r^3}\cos\theta
\end{equation}

Componente transversal:
\begin{equation}
    E_\theta = -\frac{1}{r}\frac{\partial\phi}{\partial\theta}
             = -\frac{1}{r}\!\left(\frac{V_0 R^2}{2r^2}\sin\theta\right)
             = -\frac{V_0 R^2}{2r^3}\sin\theta
\end{equation}

El campo eléctrico en el exterior ($r > R$) es:
\begin{equation}
    \vec{E} = \left(\frac{V_0 R}{2r^2} - \frac{V_0 R^2}{r^3}\cos\theta\right)\hat{e}_r
            - \frac{V_0 R^2}{2r^3}\sin\theta\;\hat{e}_\theta
\end{equation}

\subsection*{Nota: ¿Es necesario convertir $\sin$ a $\cos$ desde el inicio?}

\textbf{No.} Esa es precisamente la ventaja del método de ortogonalidad: el procedimiento es
sistemático y no requiere ninguna observación previa sobre la forma conveniente de la condición
de frontera.

En el método de inspección del documento original, el razonamiento depende de advertir
\emph{antes de plantear la expansión} que $\sin^2(\theta/2)$ admite una escritura útil con
$\cos\theta$.  En el método de ortogonalidad, en cambio, la relación
$\sin^2(\theta/2) = (1-u)/2$ aparece de forma automática al realizar la sustitución
$u = \cos\theta$, que es el cambio de variable estándar para todo integral de Legendre,
independientemente de la forma de la condición de frontera.  No se necesita ninguna
``astucia'' inicial: la función $(1-u)/2$ simplemente se descompone en $P_0(u)$ y $P_1(u)$
al evaluar los productos escalares \eqref{eq:int1}--\eqref{eq:int2}.

En resumen, la diferencia entre los dos métodos es de \emph{orden lógico}:
\begin{itemize}
    \item \textbf{Inspección:} detectar la forma de $\phi(R,\theta)$ $\;\Rightarrow\;$
          proponer la expansión correcta $\;\Rightarrow\;$ comparar coeficientes.
    \item \textbf{Ortogonalidad:} proponer la expansión general $\;\Rightarrow\;$
          calcular integrales con la fórmula de proyección $\;\Rightarrow\;$
          la forma de $\phi(R,\theta)$ se revela sola durante el cálculo.
\end{itemize}

\bibliographystyle{plainurl}
\bibliography{bibliografia}
\end{document}